\documentclass[11pt,a4paper]{article}

% -----------------------------------------------------------------------
% Packages
% -----------------------------------------------------------------------
\usepackage[T1]{fontenc}
\usepackage[utf8]{inputenc}
\usepackage{lmodern}
\usepackage{microtype}
\usepackage{geometry}
\geometry{margin=2.5cm}

\usepackage{amsmath,amssymb,amsthm}
\usepackage{bm}
\usepackage[hidelinks]{hyperref}
\usepackage{booktabs}
\usepackage{graphicx}
\usepackage{xcolor}
\usepackage{setspace}
\onehalfspacing

\usepackage[backend=biber,
            style=authoryear,
            sorting=nyt,
            maxcitenames=2,
            maxbibnames=99,
            giveninits=true,
            uniquename=init,
            doi=true,
            url=false]{biblatex}
\addbibresource{references.bib}

% -----------------------------------------------------------------------
% Custom macros
% -----------------------------------------------------------------------
\newcommand{\vect}[1]{\boldsymbol{#1}}
\newcommand{\R}{\mathbb{R}}

% -----------------------------------------------------------------------
% Meta-data
% -----------------------------------------------------------------------
\title{\textbf{Topic Modeling from Latent Semantic Analysis to
  Large Language Models:\\
  A Critical Review of Methods and Evaluation Practice}}

\author{%
  \textit{[Author name(s) to be supplied at submission]}\\[4pt]
  \textit{[Institutional affiliation(s)]}
}

\date{\today}

% -----------------------------------------------------------------------
\begin{document}
\maketitle
\begin{abstract}
Topic modeling has grown from a compact statistical idea—decomposing
a document collection into latent thematic components—into an active
cross-disciplinary methodology spanning statistics, machine learning,
and natural language processing.
This paper critically surveys the trajectory from Latent Semantic
Analysis through probabilistic models (pLSA, LDA and its many
extensions) to neural variational autoencoders, contextualised
embedding-based methods, and the most recent wave of
large-language-model (LLM) assisted topic discovery.
We devote particular attention to evaluation, arguing that the field's
most persistent and underappreciated problem is not model expressiveness
but measurement: coherence metrics widely adopted as proxies for
topic quality are themselves unreliable, and held-out perplexity
correlates poorly with human interpretability.
We review stability, reproducibility, software tooling, multilingual
extensions, applications, and emerging ethical concerns.
Throughout, we identify where methodological progress has been genuine
and where it rests on shaky evaluation foundations, and we outline
open problems that deserve sustained future attention.
\end{abstract}

\tableofcontents
\newpage

% =======================================================================
\section{Introduction}
\label{sec:introduction}
% =======================================================================

A recurring challenge in text analysis is the need to summarise and
navigate large document collections without exhaustive manual reading.
Topic modeling addresses this challenge by positing that documents are
mixtures of latent themes, and that these themes can be recovered from
word co-occurrence statistics alone.
The appeal is immediate: given a corpus of, say, scientific abstracts or
social-media posts, a topic model returns a human-readable vocabulary
for the collection's major subjects, together with document-level
scores indicating how strongly each subject is represented.

The field's intellectual history spans three decades.
\textcite{deerwester1990indexing} introduced Latent Semantic Analysis
(LSA), a linear-algebraic technique that exposed latent structure via
singular value decomposition (SVD).
\textcite{hofmann1999probabilistic} recast the problem probabilistically
with pLSA, introducing a proper generative story.
\textcite{blei2003latent} then added a Bayesian prior over document
mixtures, yielding Latent Dirichlet Allocation (LDA)---still, more than
two decades later, the most widely cited topic model.
Subsequent years brought a rich family of LDA extensions addressing
temporal dynamics \parencite{blei2006dynamic}, topic correlations
\parencite{blei2005correlated}, supervision \parencite{blei2007supervised},
hierarchical structure \parencite{teh2006hierarchical}, short text
\parencite{yin2014dirichlet}, and multiple languages
\parencite{mimno2009polylingual}.
The neural wave, beginning around 2016--2017, replaced conjugate
Bayesian inference with variational autoencoders
\parencite{miao2016neural,srivastava2017autoencoding}, enabling more
flexible parameterisation and scalability.
Embedding-based models followed, grounding topics directly in
pre-trained semantic spaces \parencite{dieng2020topic,sia2020tired,
bianchi2021pretraining,grootendorst2022bertopic}.
Most recently, large language models have been recruited as topic
discovery engines, with \textcite{pham2024topicgpt} demonstrating that
prompting an LLM with document samples can rival or exceed traditional
models on human-judged coherence.

Despite this proliferation of methods, the evaluation landscape has
lagged behind.
The standard metrics---held-out perplexity, automated coherence
scores---have been shown to be unreliable proxies for what practitioners
actually want from topic models: interpretable, stable, and actionable
thematic summaries.
The human-in-the-loop experiments of \textcite{chang2009reading}
established early on that lower perplexity does not imply better human
interpretability.
Subsequent work on coherence metrics
\parencite{newman2010automatic,lau2014machine,roder2015exploring}
brought automation to evaluation, but later studies showed these metrics
can be gamed and that their correlation with human judgment is weaker
than assumed \parencite{hoyle2021automated,doogan2021twaddle}.
Neural topic models raise further concerns: \textcite{hoyle2022broken}
argue that many neural models are ``broken'' in the sense that their
topic word lists are less meaningful than those of classical LDA despite
superior coherence scores.

This review makes three primary contributions.
First, it provides a structured synthesis of method families that
reflects both the logical and historical relationships among them,
rather than merely cataloguing papers.
Second, it gives evaluation practice a central and critical role,
tracing the gap between what is being measured and what practitioners
need.
Third, it discusses practical dimensions---tooling, reproducibility,
multilingual applicability, energy costs, and applications---that are
often treated as afterthoughts in methodology-focused surveys.

\medskip
\noindent\textbf{Scope.}\quad
This review covers peer-reviewed and directly verified work from 1990
through 2024.
We include one arXiv preprint (BERTopic; \cite{grootendorst2022bertopic})
because of its exceptional adoption in applied research, explicitly
noting its unreviewed status where relevant.
We do not attempt to cover every downstream application domain; instead
we use representative case studies to illustrate the breadth of
deployment.

% =======================================================================
\section{Foundations}
\label{sec:foundations}
% =======================================================================

\subsection{Latent Semantic Analysis}
\label{sec:lsa}

Latent Semantic Analysis \parencite{deerwester1990indexing} begins with a
term--document matrix $\mathbf{X} \in \R^{V \times D}$, where $V$ is the
vocabulary size and $D$ is the number of documents.
Truncated SVD decomposes $\mathbf{X} \approx \mathbf{U}_{k}
\mathbf{\Sigma}_{k} \mathbf{V}_{k}^{\top}$, projecting both terms and
documents into a shared $k$-dimensional semantic space.
The rank-$k$ approximation captures the dominant variance in word use,
smoothing over synonymy and polysemy in a way that raw term matching
cannot.

LSA achieves impressive results in information retrieval and
document similarity tasks, but it has well-known limitations.
The latent dimensions are not constrained to be non-negative or
interpretable; they are algebraic artefacts of the SVD rather than
meaningful topics.
Negative entries in factor matrices make direct interpretation
difficult, as a ``negative loading'' on a dimension has no
straightforward semantic gloss.
Furthermore, LSA provides no probabilistic account of document
generation, making it impossible to assign principled probability
estimates to held-out data or to perform likelihood-based model
selection.
The choice of rank $k$ is also heuristic: practitioners typically try
values in the range 50--500 and rely on downstream task performance or
qualitative inspection to guide selection.

Another limitation is that LSA treats all terms as contributing
equally to the latent space unless TF-IDF weighting is applied prior
to decomposition.
Function words and highly frequent terms can dominate the decomposition
unless they are removed or down-weighted.
These preprocessing decisions can substantially change the latent
dimensions recovered, and they interact with the choice of $k$ in
ways that are difficult to predict.
Despite these limitations, LSA demonstrated conclusively that latent
semantic structure could be extracted from large corpora by purely
statistical means, without any hand-crafted ontology or lexical
resource, and this insight directly motivated the probabilistic
reformulations that followed.

\subsection{Probabilistic Latent Semantic Indexing}
\label{sec:plsa}

\textcite{hofmann1999probabilistic} addressed the probabilistic
deficiency of LSA by introducing pLSA (also written PLSI).
The generative story posits a latent topic variable $z \in
\{1,\ldots,K\}$: a word $w$ in document $d$ is generated by first
sampling a topic $z \sim P(z|d)$ and then sampling the word from the
topic-specific distribution $P(w|z)$.
The parameters $P(z|d)$ and $P(w|z)$ are estimated by expectation
maximisation (EM), giving a clean probabilistic interpretation.

pLSA improves on LSA in information retrieval benchmarks and provides
a likelihood score, but it is not a fully generative model of new
documents: the per-document parameters $P(z|d)$ are fixed to the
training corpus and do not generalise to unseen texts.
This is the \emph{train/test disconnect}---the model cannot assign a
probability to a document outside the training set without additional
folding-in steps.
The number of parameters grows linearly with the corpus size,
$\mathcal{O}(D \cdot K + V \cdot K)$, raising overfitting concerns
on small corpora and making storage and serialisation unwieldy for
very large collections.
Regularisation strategies such as tempered EM were proposed to
address overfitting, but they do not resolve the fundamental
generalisation problem.

The train/test disconnect in pLSA is not merely a theoretical concern.
It means that two independently trained pLSA models on the same corpus
are not directly comparable: there is no natural alignment of topic
indices between runs, and no principled way to apply a trained model
to compute the likelihood of new data.
These properties make pLSA unsuitable for the growing class of
applications that require either online processing of new documents or
quantitative model comparison via held-out likelihood.

\subsection{Latent Dirichlet Allocation}
\label{sec:lda}

\textcite{blei2003latent} resolved the generalisation problem by placing
Dirichlet priors over both the per-document topic distributions
$\vect{\theta}_{d}$ and the per-topic word distributions
$\vect{\beta}_{k}$.
The full generative process for LDA is:
\begin{enumerate}
  \item For each topic $k = 1, \ldots, K$: draw
    $\vect{\beta}_{k} \sim \mathrm{Dir}(\vect{\eta})$.
  \item For each document $d = 1, \ldots, D$:
    \begin{enumerate}
      \item Draw $\vect{\theta}_{d} \sim \mathrm{Dir}(\vect{\alpha})$.
      \item For each word position $n = 1, \ldots, N_{d}$:
        \begin{enumerate}
          \item Draw $z_{d,n} \sim \mathrm{Categorical}(\vect{\theta}_{d})$.
          \item Draw $w_{d,n} \sim \mathrm{Categorical}(\vect{\beta}_{z_{d,n}})$.
        \end{enumerate}
    \end{enumerate}
\end{enumerate}
The Dirichlet prior over $\vect{\theta}_{d}$ enables the model to
assign probabilities to new documents via the predictive distribution
$p(\mathbf{w}|\vect{\alpha},\vect{\beta})$, solving the train/test
disconnect.

Posterior inference in LDA is intractable in closed form because the
document-topic distribution $\vect{\theta}_{d}$ and the topic
assignments $z_{d,n}$ are coupled through the observation $w_{d,n}$.
The original paper \parencite{blei2003latent} employed mean-field
variational EM, which introduces a fully-factored variational distribution
$q(\vect{\theta},\mathbf{z})$ and optimises its parameters to minimise
the KL divergence from the true posterior.
\textcite{griffiths2004finding} introduced collapsed Gibbs sampling as
an alternative: by integrating out $\vect{\theta}_{d}$ and $\vect{\beta}_{k}$
analytically, the sampler operates only over the topic assignment
variables $z_{d,n}$, yielding simpler update equations and stronger
mixing in practice.
Collapsed Gibbs sampling became the dominant inference strategy for LDA
and remains the reference implementation in tools such as MALLET.
More recent work explored stochastic variational inference
\parencite{blei2012probabilistic}, enabling online processing of
streaming corpora that do not fit in memory.

\textcite{blei2012probabilistic} provides an accessible tutorial
covering LDA and its probabilistic semantics for a broad audience;
it remains the most widely cited entry point for practitioners new to
probabilistic topic models.

LDA treats the $K$ topics as a priori independent and exchangeable,
and it assumes the bag-of-words representation: word order and syntax
are discarded.
These simplifications are strong but not always damaging;
for many corpus-level analyses the loss of sequential information
matters less than the ability to summarise tens of thousands of
documents cheaply and interpretably.
The hyperparameter $\vect{\alpha}$ controls topic sparsity in
documents: small symmetric $\alpha$ encourages documents to be
concentrated on few topics, while large $\alpha$ pushes documents
toward uniform mixtures.
Similarly, $\vect{\eta}$ controls sparsity in topic-word distributions.
Empirically, asymmetric priors---particularly a peaked prior over
the topic distribution---often yield more stable and interpretable
topics than the symmetric defaults, a finding with practical
implications for deployment \parencite{koltcov2014lda}.

The appropriate number of topics $K$ must be chosen by the analyst,
typically via grid search over held-out likelihood or coherence, a
pragmatic but imperfect procedure.
No single criterion for choosing $K$ has proven universally reliable:
the value that maximises held-out perplexity often differs from the
value that maximises human-judged coherence, and both may differ from
the value that is most useful for a downstream task.
This $K$-selection problem has motivated nonparametric alternatives
(Section~\ref{sec:hdp}) but remains unsolved even for practitioners
who prefer parametric models.

% =======================================================================
\section{Extensions to Latent Dirichlet Allocation}
\label{sec:lda-extensions}
% =======================================================================

The success of LDA prompted a wave of extensions addressing its
structural assumptions.
We organise these by the dimension of variation they introduce.

\subsection{Topic Correlation: the Correlated Topic Model}
\label{sec:ctm}

LDA's Dirichlet prior imposes negative correlations between topic
proportions: documents heavy on one topic tend to be lighter on
others, a consequence of the Dirichlet's simplicial geometry.
\textcite{blei2005correlated} replaced the Dirichlet with a logistic
normal distribution, yielding the Correlated Topic Model (CTM).
The logistic normal can capture positive topic correlations
(\emph{politics} and \emph{economics} may co-occur) as well as
negative ones.
The trade-off is computational: the logistic normal does not have
Dirichlet-conjugate conditional distributions, so Gibbs sampling is
not straightforward and variational inference requires a further
approximation layer.
On scientific corpora the CTM reliably discovers meaningful topical
correlations, but the gains over LDA in downstream retrieval tasks are
often modest.

\subsection{Temporal Dynamics: Dynamic Topic Models}
\label{sec:dtm}

Document collections often span time, and topics evolve.
\textcite{blei2006dynamic} introduced Dynamic Topic Models (DTM) by
placing a linear Gaussian state-space model over sequential topic
word distributions: the topic $\vect{\beta}_{k,t}$ at time slice $t$
is a Gaussian perturbation of the topic at $t-1$.
This allows a historian to trace, for instance, how the vocabulary
associated with the topic \emph{artificial intelligence} shifted from
the 1950s through the 2010s.
The DTM has been applied to scientific literature, news archives, and
congressional records.
The principal limitation is scalability: exact inference is intractable,
and the variational approximation is substantially more involved than
standard LDA inference.

\subsection{Nonparametric Approaches: Hierarchical Dirichlet Processes}
\label{sec:hdp}

Both LDA and DTM require the analyst to specify $K$.
\textcite{teh2006hierarchical} proposed the Hierarchical Dirichlet
Process (HDP), a nonparametric Bayesian model in which $K$ is inferred
from data.
The HDP places a Dirichlet Process prior over a global measure that
is shared across documents; each document's topic distribution is
a draw from a document-level DP that is centred on the global measure.
In practice, the effective number of topics grows logarithmically
with corpus size, providing automatic complexity regularisation.
HDP inference is considerably more complex than LDA inference and
convergence can be slow; moreover, the number of inferred topics is
sensitive to hyperparameter settings, so the claim of fully automatic
topic number selection should be tempered.

\subsection{Supervised Topic Models}
\label{sec:slda}

\textcite{blei2007supervised} introduced Supervised LDA (sLDA) by
adding a response variable $y_{d}$---a rating, label, or continuous
outcome---that is modelled as a function of the document's topic
proportions $\vect{\theta}_{d}$.
Supervised inference encourages the model to discover topics that are
predictive of the outcome, yielding more task-relevant thematic
summaries than unsupervised LDA when labelled data are available.
A complementary approach, Labeled LDA \parencite{ramage2009labeled},
restricts each document's topic prior to the subset of topics
corresponding to its observed labels, providing a finer-grained credit
attribution mechanism in multi-label settings such as web-page
categorisation.

\subsection{Short-Text and Sparse Settings}
\label{sec:shorttext}

Standard LDA assumes documents contain enough words for reliable
estimation of $\vect{\theta}_{d}$; tweets, titles, and product reviews
often contain fewer than thirty tokens.
\textcite{yin2014dirichlet} addressed this with the Dirichlet
Multinomial Mixture (DMM) model, which assigns each short document to
a single topic rather than a mixture, reducing the effective number of
latent variables per document.
On several benchmark short-text clustering tasks, DMM outperforms LDA
by a wide margin while remaining interpretable.
The single-topic assumption is strong and fails for longer, heterogeneous
documents, so DMM is best understood as a specialised rather than a
universal replacement for LDA.

\subsection{Multilingual Topic Models}
\label{sec:multilingual}

Standard LDA operates over a fixed vocabulary in a single language.
\textcite{mimno2009polylingual} extended LDA to aligned multilingual
corpora---parallel or comparable documents in multiple languages---by
assuming that aligned documents share the same topic distribution
$\vect{\theta}_{d}$ but each language maintains its own topic-word
distribution.
This allows topics to be inferred jointly from, for example, English
and Chinese text, with the semantic coherence of the shared topic
serving as a cross-lingual anchor.
More recent work \parencite{bianchi2021crosslingual} incorporates
contextualised multilingual embeddings (e.g.\ mBERT, XLM-R) as
document representations, enabling zero-shot topic inference in
languages for which no parallel training data are available.
We return to cross-lingual approaches in Section~\ref{sec:contextualized}.

% =======================================================================
\section{Neural and Embedding-Based Topic Models}
\label{sec:neural}
% =======================================================================

The deep-learning revolution brought two major shifts to topic modeling:
(i) replacing conjugate Bayesian inference with amortised variational
autoencoders, and (ii) replacing bag-of-words representations with
dense semantic embeddings from pre-trained language models.

\subsection{Neural Variational Document Models}
\label{sec:nvdm}

\textcite{miao2016neural} introduced the Neural Variational Document
Model (NVDM), which encodes a document's term frequency vector through
a neural network into a Gaussian latent code, then decodes that code
into word probabilities.
The encoder--decoder architecture is trained end-to-end via the
Evidence Lower Bound (ELBO) of variational inference.
NVDM was an important proof of concept: it demonstrated that topic-like
representations could be learned without explicit Dirichlet priors or
Gibbs sampling.
However, the Gaussian latent space does not enforce the simplex
structure of topic proportions, and early versions produced less
coherent topics than well-tuned LDA.

\textcite{srivastava2017autoencoding} addressed this with the
Autoencoded Variational Inference for Topic Models (AVITM) framework,
which approximates the Dirichlet distribution using a logistic-normal
auxiliary variable---a technique sometimes called the Laplace
approximation to the Dirichlet.
AVITM introduced the key engineering insight of \emph{reparameterisation}
through a logistic-normal: by sampling from the Gaussian and passing
through a softmax, the decoder receives simplex-valued topic proportions
that closely mirror the Dirichlet.
Comparisons with LDA on standard benchmarks showed competitive or
superior coherence and dramatically faster training on GPU hardware.
The correct OpenReview identifier for this work is \texttt{BybtVK9lg}
(not the erroneous \texttt{HkwzHslxg} cited in some secondary sources).

\textcite{zhao2021neural} propose a further departure, replacing the
ELBO with an Optimal Transport (OT) objective.
Their Neural Sinkhornian Topic Model directly minimises the Wasserstein
distance between document word distributions and reconstructed topic
mixtures, learning topic embeddings jointly in the OT cost matrix.
This formulation is theoretically principled and shows competitive
results on both coherence and topic diversity metrics, though the OT
computation introduces additional hyperparameter sensitivity.

\subsection{Embedding-Based Models: Topics as Vectors}
\label{sec:etm}

\textcite{dieng2020topic} introduced the Embedded Topic Model (ETM),
which places topics and vocabulary items in the same embedding space
$\R^{L}$.
The word distribution for topic $k$ is the softmax of the dot products
between the topic embedding $\vect{\alpha}_{k} \in \R^{L}$ and all word
embeddings $\vect{\rho}_{v} \in \R^{L}$:
\[
  p(w = v \mid z = k) \propto \exp\bigl(\vect{\rho}_{v}^{\top}
  \vect{\alpha}_{k}\bigr).
\]
The topic and word embeddings can be initialised from pre-trained
word2vec or GloVe vectors, injecting rich semantic knowledge into the
model without requiring a separate fine-tuning step.
ETM achieves better coherence on large vocabularies than standard LDA
and scales gracefully to corpora with $10^{5}$--$10^{6}$ documents.

\textcite{sia2020tired} take a different perspective, arguing that
off-the-shelf clustering of pre-trained word embeddings produces topics
that are fast, competitive in coherence, and transparent---without any
probabilistic modelling at all.
Their empirical comparison raises an important question: if simple
$k$-means over word vectors achieves similar quality to sophisticated
models, what is the value-add of the additional modelling complexity?
The honest answer is that it depends on the task.
Probabilistic models offer principled uncertainty quantification,
document-level topic scores, and a generative story that supports
tasks like document retrieval and completion.
Clustering methods offer speed and simplicity.
Both have a role, and the choice should be driven by the downstream
use case.

\subsection{Contextualised and Cross-Lingual Models}
\label{sec:contextualized}

Pre-trained transformer models such as BERT encode context-sensitive
word representations, and they can be used to produce document-level
embeddings that are richer than bag-of-words vectors.
\textcite{bianchi2021pretraining} demonstrate that substituting BERT
embeddings for TF-IDF vectors in a neural topic model substantially
improves topic coherence on standard benchmarks.
They argue that pre-training effectively provides a strong semantic
inductive bias that the model would otherwise need to learn from the
target corpus alone---a significant benefit on small or domain-specific
corpora.

The cross-lingual extension \parencite{bianchi2021crosslingual} uses
multilingual pre-trained encoders to project documents in different
languages into a shared representation space.
Topics are then inferred in this shared space, and zero-shot transfer
allows topic models trained on English to be directly applied to
languages not seen during training.
Empirical results on English--Italian and English--German corpora show
substantially better cross-lingual topic alignment than earlier
dictionary-based or translation-based approaches.

\textcite{grootendorst2022bertopic} presents BERTopic, an applied
framework that clusters BERT sentence embeddings using HDBSCAN and then
characterises each cluster with a class-based TF-IDF score.
BERTopic is a preprint (arXiv:2203.05794) that has not undergone
formal peer review, yet it has been adopted in numerous applied
projects due to its ease of use and integration with the
\texttt{sentence-transformers} ecosystem.
Its coherence scores on standard benchmarks are competitive with
AVITM and ETM, but independent evaluations have noted sensitivity
to HDBSCAN hyperparameters and a tendency to produce one very large
miscellaneous cluster when the underlying data distribution is diffuse.

% =======================================================================
\section{LLM-Assisted Topic Discovery}
\label{sec:llm}
% =======================================================================

The most recent paradigm shift recruits large language models not as
representation providers but as direct topic-discovery agents.
\textcite{pham2024topicgpt} propose TopicGPT, which prompts an
instruction-tuned LLM (GPT-4 in their experiments) with a small
sample of documents and asks it to produce a set of topic labels
together with evidence passages.
The framework supports iterative refinement: the analyst can provide
feedback and the model revises its topic assignments accordingly.
On human evaluation benchmarks, TopicGPT-produced topics are judged
more coherent and informative than LDA or AVITM baselines in three out
of four domains tested.

The attractions of LLM-based approaches are considerable: they can
leverage world knowledge beyond the target corpus, produce
naturally readable labels without post-hoc aggregation, and adapt
interactively to analyst needs.
The framework supports iterative refinement at both the granularity
level (splitting fine-grained topics or merging redundant ones) and
the label level (instructing the model to use domain-specific
terminology), a degree of human-in-the-loop flexibility that
probabilistic models cannot easily provide.

However, several concerns merit attention.
First, computational cost is high: querying a proprietary LLM API for
every document at scale is expensive both financially and in energy
terms \parencite{strubell2019energy}, whereas LDA inference on a
million-document corpus can run on a single CPU in hours.
Second, LLM outputs are not probabilistic in the topic-model sense:
there are no well-defined $\vect{\theta}_{d}$ vectors that can be
fed into downstream statistical models.
This means that LLM-discovered topics are not directly usable as
features for classification or retrieval without an additional
assignment step.
Third, reproducibility is constrained by model versioning and API
availability; a study using GPT-4 today cannot guarantee the same
outputs next year, making longitudinal comparisons and exact
replication impossible in principle.
Fourth, the evaluation of LLM-topic quality itself reintroduces the
problems described in Section~\ref{sec:evaluation}: human raters
have been shown to favour fluent, familiar language irrespective of
the underlying analytical validity of the topic groupings.
Topics produced by an LLM may score highly on human coherence tasks
not because they are analytically meaningful but because they are
grammatically polished and semantically familiar.

Fifth, the boundaries between what the LLM has learned from
pre-training data and what it has discovered from the target corpus
are opaque.
When TopicGPT produces a topic labelled ``climate policy,'' it is
unclear how much of that topic reflects genuine signals in the target
documents versus the LLM's prior associations between related
concepts.
This opacity complicates the use of LLM-based topic models in
interpretive social-science research, where the distinction between
what the data say and what the analyst's tools assume is critical.

A theoretically sound integration of LLM representations into
probabilistic topic models remains an open research problem.
Current approaches are largely pipeline-based---using LLM embeddings
as inputs to a conventional topic model---rather than architecturally
unified.

% =======================================================================
\section{Evaluation: Metrics, Pitfalls, and Recommendations}
\label{sec:evaluation}
% =======================================================================

Evaluation is arguably the most intellectually fraught aspect of topic
modeling.
We organise the discussion around four families of metrics: (i)
held-out likelihood, (ii) intrinsic coherence, (iii) human evaluation,
and (iv) extrinsic/downstream utility.

\subsection{Held-Out Likelihood and Perplexity}
\label{sec:eval-perplexity}

Held-out perplexity is the canonical probabilistic model-selection
criterion.
\textcite{wallach2009evaluation} give an authoritative treatment of how
to compute held-out log-likelihood for LDA and its extensions, noting
that naive approaches (computing the marginal likelihood using only the
training-set parameters) are biased and that correct estimation requires
either left-to-right sequential Monte Carlo, harmonic mean estimation
with appropriate care, or AIS.
Their experiments show that different estimators can give substantially
different rankings of models, underscoring the need for careful
implementation.

The fundamental limitation of perplexity as an evaluation criterion was
demonstrated by \textcite{chang2009reading} through a series of human
experiments.
Their \emph{word intrusion} task presents a human judge with the top
words of a topic plus one intruder word randomly drawn from another
topic; a high-quality topic should make the intruder easy to detect.
Their \emph{topic intrusion} task similarly evaluates whether humans
can identify which topic was \emph{not} used to generate a document.
Across multiple models and datasets, models with lower perplexity often
performed \emph{worse} on these human-centered tasks, a striking finding
that motivated decades of subsequent evaluation research.

\subsection{Automatic Coherence Metrics}
\label{sec:eval-coherence}

\textcite{newman2010automatic} proposed one of the first automated
coherence metrics, based on pointwise mutual information (PMI) between
pairs of top topic words computed over a large reference corpus.
Topics whose top words co-occur more than chance exhibit higher PMI,
and such topics tend to be rated more coherent by human judges.
\textcite{lau2014machine} extended this work by fitting regression
models to human coherence ratings, finding that normalised PMI (NPMI)
computed over a Wikipedia reference corpus correlates most strongly
with human judgments.
\textcite{roder2015exploring} systematically catalogued a large family
of coherence measures (UCI, UMass, NPMI, $C_{v}$, $C_{p}$, etc.) and
found that $C_{v}$ and $C_{NPMI}$ provide the best approximation of
human judgment across datasets, while UMass---the simplest and most
computationally efficient measure---correlates least well.

Despite this progress, serious doubts about automated coherence metrics
have accumulated.
\textcite{mimno2011optimizing} introduced the UMass coherence measure
as a fast, corpus-internal alternative, but also showed that directly
optimising it during inference improves topic quality---raising the
spectre of optimising the surrogate rather than the underlying
objective.
\textcite{hoyle2021automated} conduct a large-scale empirical audit
and show that standard coherence metrics such as NPMI and $C_{v}$
are poorly calibrated: models can achieve high coherence scores
while producing topics that human evaluators judge as nonsensical.
The correlation between automatic and human scores degrades further
when the reference corpus for PMI computation differs from the
target corpus---a common mismatch in practice.
\textcite{doogan2021twaddle} reach similar conclusions through a
complementary analysis, arguing that topic coherence as currently
measured is largely a measure of corpus-frequency effects and generic
collocations rather than genuine semantic interpretability.

The practical implication is that researchers should report at least
one human evaluation component alongside any automated coherence
score.
Neither metric alone is sufficient.

\subsection{Stability and Reproducibility}
\label{sec:eval-stability}

A topic model is stable if re-running the same algorithm on the same
corpus produces essentially the same topics.
\textcite{koltcov2014lda} document that LDA with collapsed Gibbs
sampling can exhibit substantial run-to-run variation in topic
assignments, particularly for the default symmetric Dirichlet prior
settings.
They find that topic stability improves with more Gibbs iterations and
with asymmetric priors, but remains an underreported and underappreciated
source of variance in empirical comparisons.
Stability is especially relevant for neural models, where stochastic
gradient descent and random weight initialisation add further sources
of variation.
Standard practice should include reporting variance across multiple
random seeds, but this is often omitted in papers that report a single
run.

\subsection{Extrinsic Evaluation and Tooling}
\label{sec:eval-extrinsic}

Intrinsic metrics measure topic quality in isolation; extrinsic metrics
evaluate topics in the context of a downstream task---classification,
retrieval, recommendation, or social-science analysis.
Extrinsic evaluation is more directly aligned with user goals but
requires task-specific annotation and makes comparison across papers
difficult.

\textcite{terragni2021octis} provide OCTIS (Optimizing and Comparing
Topic models Is Simple), an open-source benchmarking framework that
standardises training, evaluation, and hyperparameter search for a
wide range of topic models including LDA, AVITM, CTM, ETM, NMF, and
others.
OCTIS reports a battery of coherence, diversity, and classification
metrics on a common set of preprocessed benchmark corpora (20~Newsgroups,
BBC News, M10, and others), enabling more systematic comparison than
ad-hoc per-paper evaluations.
Topic diversity---the proportion of unique words in the top-$N$ words
across all topics---is included as a complementary metric to coherence,
addressing the concern that models may produce high-coherence but highly
redundant topics.
OCTIS has become a de-facto standard for neural topic model evaluation
and is actively maintained.

The findings from OCTIS and related benchmarking studies paint a
nuanced picture.
Classical LDA remains a strong baseline: it is hard to beat
consistently on both coherence and diversity simultaneously.
Neural models show higher variance across runs and datasets.
Hyperparameter choices, particularly the number of topics and
the encoder architecture, have a larger effect on outcome than the
choice of model family, a finding that has implications for how papers
should report and ablate their results.

An important caveat about topic diversity is worth highlighting.
Topic diversity---computed as the fraction of unique top-$N$ words
across all $K$ topics---is a useful complement to coherence, because
a model can achieve high per-topic coherence while producing many
nearly-identical topics that collectively cover only a narrow slice of
the corpus's thematic range.
However, diversity is not a goal in itself: a corpus genuinely dominated
by a single theme should produce low diversity, and penalising such
a model is unfair.
As with all single-metric evaluations, diversity scores should be
interpreted in light of the corpus characteristics and the downstream
use case.

\subsection{Human Evaluation Best Practices}
\label{sec:eval-human}

\textcite{chang2009reading} established word intrusion and topic
intrusion as the reference protocol for human evaluation.
Both tasks can be implemented as crowd-sourced studies with reasonable
inter-annotator agreement ($\kappa > 0.6$ in well-controlled settings).
Their finding that perplexity anticorrelates with human interpretability
should be treated as a standing reminder: when evaluating a new model,
always include at least one human-centered task.

The word intrusion task is particularly informative about topic purity:
a human judge who can consistently identify the randomly injected
``intruder'' word confirms that the remaining top words form a
coherent, distinctive cluster.
Conversely, a topic in which the intruder is invisible to judges is
one whose words are semantically diffuse or generic.

An important nuance is that coherence and usefulness are not the same.
Topics judged coherent by crowdworkers may still be uninformative from
an analytical standpoint if they capture superficial lexical
associations rather than substantive themes.
Consider a topic dominated by grammatical function words or generic
academic vocabulary: it will score poorly on PMI-based coherence but
may still appear coherent to a non-specialist crowdworker.
Conversely, a highly domain-specific topic may appear incoherent to
a general crowdworker who lacks the domain background to recognise
the semantic relationships.
Designing evaluation protocols that test \emph{analytical utility}
rather than mere \emph{word association} remains an open challenge.

The relationship between automatic coherence scores and task utility
has been studied most directly in classification and retrieval
settings.
OCTIS reports classification accuracy (using topic vectors as features)
alongside coherence, allowing researchers to observe cases where the
two metrics diverge.
Such divergences are informative: they reveal that a model producing
high-coherence topics is not necessarily producing better document
representations for downstream prediction, reinforcing the case for
multi-metric evaluation.

% =======================================================================
\section{Comparative Analysis}
\label{sec:comparison}
% =======================================================================

Table~\ref{tab:model-comparison} summarises key properties of the
major model families reviewed above.
Across every dimension, no single model dominates all others.

\begin{table}[ht]
\centering
\caption{Comparative summary of topic model families. ``Inf.'' =
  inference algorithm; ``Prob.\ doc.'' = probabilistic document
  representation; ``Scalability'' = approximate assessment based on
  reported corpus sizes in original papers.}
\label{tab:model-comparison}
\small
\begin{tabular}{@{}lllll@{}}
\toprule
Model & Inf. & Prob.\ doc. & Interpretability & Scalability \\
\midrule
LSA                        & SVD              & No  & Low   & High \\
pLSA                       & EM               & Yes & Med.  & Med. \\
LDA                        & Gibbs/Var.       & Yes & High  & High \\
CTM                        & Var.             & Yes & Med.  & Med. \\
DTM                        & Var.             & Yes & High  & Low  \\
HDP                        & Gibbs            & Yes & High  & Med. \\
sLDA                       & Var.             & Yes & High  & Med. \\
DMM                        & Gibbs            & Yes & High  & High \\
NVDM / AVITM               & VAE              & Yes & Med.  & High \\
ETM                        & VAE              & Yes & High  & High \\
BERTopic                   & Clustering       & No  & Med.  & High \\
TopicGPT                   & LLM prompting    & No  & High* & Low  \\
\bottomrule
\multicolumn{5}{l}{* Human-judged; not probabilistic.}
\end{tabular}
\end{table}

Several cross-cutting observations are worth highlighting:

\paragraph{Flexibility vs.\ interpretability.}
More flexible models (DTM, CTM, neural VAEs) gain representational
power at the cost of interpretability and stability.
The extra parameters can be estimated reliably only with large corpora;
on small to medium-sized collections (fewer than $\sim$10,000
documents), simple LDA often produces more consistent topics.
Practitioners working with limited data are better served by a
well-tuned LDA with asymmetric priors than by a neural model whose
variational parameters cannot be reliably estimated from a few thousand
documents.

\paragraph{Embeddings and pre-training.}
Pre-trained embeddings provide an effective inductive bias that can
substitute for corpus size.
ETM and AVITM with pre-trained word vectors outperform their
randomly-initialised counterparts, particularly on domain-specific
corpora where the general vocabulary statistics differ from the target.
The benefit is most pronounced when the pre-training corpus is
semantically related to the target domain---medical pre-training helps
clinical topic models more than general web-text pre-training.

\paragraph{LLM methods and the probability gap.}
LLM-based approaches such as TopicGPT produce impressively readable
topics, but the absence of a probabilistic generative model means that
standard evaluation metrics (perplexity, per-document $\vect{\theta}$)
are inapplicable.
They also inherit the substantial energy cost and reproducibility
limitations of large model inference \parencite{strubell2019energy}.
Whether LLM-discovered topics are ``better'' in any absolute sense
depends entirely on what the analyst means by better: if the goal is
a readable, analyst-curated taxonomy for qualitative analysis, LLMs
are compelling; if the goal is a set of document-level numeric
representations for downstream modelling, a probabilistic topic model
is currently indispensable.

\paragraph{The inference efficiency frontier.}
Gibbs sampling for LDA scales well to millions of documents via
sparse implementations, but is serial and slow to parallelize.
VAE-based models benefit from GPU acceleration and mini-batch training,
converging orders of magnitude faster on large corpora.
BERTopic's pipeline---embedding, dimensionality reduction (UMAP),
clustering (HDBSCAN), topic characterisation (c-TF-IDF)---is
parallelisable and fast, but its quality is sensitive to a three-way
hyperparameter interaction (embedding model, UMAP parameters, HDBSCAN
parameters) that is difficult to tune systematically.

\paragraph{Evaluation as a first-class concern.}
Across all model families, evaluation practice has consistently lagged
behind model development.
Papers routinely report a single automated coherence score without
stability analysis, human evaluation, or extrinsic benchmarks.
The evidence from \textcite{hoyle2021automated},
\textcite{doogan2021twaddle}, and \textcite{hoyle2022broken} suggests
that this gap is not merely a reporting convenience but can lead to
substantively misleading conclusions about model quality.
A model that achieves the highest $C_{v}$ coherence on one dataset may
rank last on human-judged interpretability on a related dataset;
without reporting both, neither researcher nor practitioner can know
which conclusion to trust.

% =======================================================================
\section{Applications}
\label{sec:applications}
% =======================================================================

Topic models have been applied across a wide range of domains.
We highlight a few illustrative cases to convey the breadth of
deployment and the practical evaluation considerations that arise.

\paragraph{Social and political science.}
\textcite{bail2016combining} demonstrate an early integration of topic
modeling with network analysis to study how advocacy organisations
stimulate political conversation on social media.
Topic models identified the thematic content of organisational posts,
while network centrality measures tracked information diffusion.
The combination revealed patterns of framing and agenda-setting
invisible to either method alone.
This study exemplifies a broader pattern in computational social science:
topic models serve as a dimensionality-reduction and coding step,
translating large text corpora into variables that can be entered into
conventional statistical models.
For this use case, the most important model property is not raw
coherence but \emph{construct validity}: do the topics correspond
to theoretically meaningful categories in the substantive domain?
No automated metric currently assesses construct validity directly,
which is why social-science applications frequently supplement
automated evaluation with expert annotation or qualitative
topic inspection.

\paragraph{Scientometrics and intellectual history.}
\textcite{griffiths2004finding} apply LDA to the PNAS corpus and show
that topic trajectories over decades align closely with the known
history of scientific discovery, providing a computationally grounded
approach to the sociology of science.
DTM has been similarly applied to trace the emergence and decline of
research themes in physics, computer science, and medicine.
These temporal applications highlight a requirement often absent in
benchmark evaluations: the model must produce topics that are stable
and interpretable not just at a single time point but across a
temporal sequence, so that year-over-year topic trajectories are
analytically meaningful rather than artefacts of random variation.

\paragraph{Information retrieval and recommendation.}
LDA's document-level topic vectors $\vect{\theta}_{d}$ can be used as
features in retrieval models, yielding document representations that
generalise beyond exact term match.
Retrieval using latent topic representations is particularly valuable
for queries that use different vocabulary than the target documents
(the synonymy problem that originally motivated LSA).
The sLDA framework \parencite{blei2007supervised} directly optimises
topic discovery for prediction of review ratings or document labels,
making it well-suited to recommendation and classification tasks where
a response variable is available.
When evaluation is conducted on a retrieval or classification task,
the superiority of LDA over bag-of-words representations depends
strongly on corpus size and the degree of vocabulary mismatch between
queries and documents; on large corpora with dense vocabulary overlap,
TF-IDF often remains competitive.

\paragraph{Public health and news analysis.}
Topic models have been applied to surveillance of social media for
emerging health crises, monitoring of media coverage of public health
events, and analysis of patient narratives in clinical records.
Short-text models such as DMM \parencite{yin2014dirichlet} are
particularly appropriate for social-media surveillance, where messages
are short and written in informal language.
In each of these settings, the primary practical requirement is not
that the model achieve the highest coherence score on a benchmark
but that it produce actionable and interpretable summaries for domain
experts---a reminder that extrinsic and human-centered evaluation
should drive model choice in applied settings rather than
performance on standardised benchmarks constructed for a different purpose.

% =======================================================================
\section{Software, Reproducibility, and Datasets}
\label{sec:software}
% =======================================================================

Reproducibility in topic modeling faces at least three challenges:
stochastic inference, sensitive hyperparameter settings, and the
absence of standardised preprocessing pipelines.
Each challenge has received some attention in the literature, but none
has been fully resolved, and their combination means that replicating
the numerical results of a typical topic-modeling paper is
substantially harder than it should be.

\paragraph{Software.}
MALLET (Machine Learning for Language Toolkit) is the most widely
used implementation of collapsed Gibbs sampling for LDA and remains a
reference baseline across the topic modeling literature.
Its efficient implementation using sparse representations of
topic-word counts enables processing of corpora with millions of
documents on a single workstation.
The Gensim library provides LDA and related models in Python with
support for online variational inference, enabling incremental training
on streaming corpora that cannot be loaded into memory at once.
Gensim also provides implementations of LSA (via SVD) and word2vec,
making it a convenient single-library environment for experiments that
span classical and neural approaches.

For neural topic models, there is no equivalent universal toolkit:
most research papers release standalone PyTorch implementations
tailored to their specific architecture, with varying degrees of
documentation and maintenance.
This fragmentation is a reproducibility concern: it is difficult to
compare a new model against a full set of baselines when each baseline
requires a separate codebase with different preprocessing assumptions.

OCTIS \parencite{terragni2021octis} directly addresses this concern by
providing standardised pipelines for multiple neural and classical
topic models, reproducible preprocessing, and a unified evaluation
harness that reports coherence, diversity, and classification metrics
on common benchmark corpora.
OCTIS has become a de-facto standard for neural topic model evaluation
and is actively maintained.
Its main limitation is coverage: at the time of writing, it does not
include LLM-based methods such as TopicGPT, reflecting the rapid pace
of development in this area.

\paragraph{Benchmark datasets.}
The 20~Newsgroups and BBC~News datasets have become de-facto standard
evaluation corpora in the topic modeling literature.
20~Newsgroups contains approximately 20,000 posts across 20 topical
categories, providing ground-truth labels that can be used for
extrinsic evaluation of topic purity.
BBC~News contains about 2,225 articles across five categories.
Both are available pre-processed through OCTIS.
The M10 dataset \parencite{terragni2021octis} is a short-text
benchmark derived from micro-posts, useful for evaluating DMM-style
models that are designed for sparse documents.

It is important to note that even these ``standard'' datasets are not
truly standardised: different papers use different tokenisation,
stop-word lists, vocabulary thresholds, and train/test splits.
\textcite{terragni2021octis} documented how preprocessing choices
affect reported coherence by several coherence-score units---enough to
reverse rankings between models.
Researchers should be explicit about which version of a dataset they
use (e.g.\ whether headers and footers are stripped from 20~Newsgroups)
and the exact preprocessing steps applied, as these choices can
substantially affect results and comparison validity.

\paragraph{Preprocessing sensitivity.}
Stop-word removal, stemming/lemmatisation, and vocabulary pruning are
preprocessing decisions with outsized influence on topic quality.
Removing too few stop words leaves topics dominated by function words;
removing too many may strip semantically relevant words.
Stemming can improve topic coherence by merging inflected forms, but it
can also conflate genuinely distinct terms (e.g.\ ``mining'' for
mineral extraction and ``mining'' in data analysis).
Vocabulary pruning---typically removing terms that appear in fewer than
$n$ documents or more than $p$\% of documents---controls sparsity and
vocabulary size, but the thresholds are rarely justified empirically.

A model that achieves high coherence on a heavily preprocessed corpus
may produce qualitatively different topics on the same corpus with
different preprocessing.
Reporting the full preprocessing pipeline is therefore as important
as reporting model hyperparameters, and benchmarks should standardise
preprocessing alongside models.
The OCTIS framework sets a useful precedent here, though extending its
preprocessing standardisation to non-English corpora remains a gap.

% =======================================================================
\section{Ethical Dimensions}
\label{sec:ethics}
% =======================================================================

\subsection{Fairness and Representation}
\label{sec:fairness}

Topic models trained on biased corpora inherit and can amplify those
biases.
A model trained predominantly on English-language, Western-authored
text will learn topics that marginalise or misrepresent other cultural
and linguistic perspectives.
Pre-trained embedding-based models (ETM, BERTopic) inherit the biases
of their upstream embeddings, which have been extensively documented
to encode racial, gender, and socioeconomic stereotypes.
Cross-lingual topic models \parencite{bianchi2021crosslingual,
mimno2009polylingual} partially address the language marginalisation
problem but do not resolve corpus-selection bias.
Practitioners should audit topic outputs for demographic and cultural
skew before deploying them in any consequential setting.

\subsection{Privacy}
\label{sec:privacy}

Topic models trained on private or sensitive text (medical records,
private messages) can expose individual-level information through
their topic-word distributions, particularly in small corpora where
unusual vocabulary items can identify specific authors.
The standard LDA generative model does not include any privacy
guarantee; differential privacy mechanisms have been proposed for
Gibbs-sampled LDA but are not yet standard practice.
Researchers working with sensitive corpora should consult the
differential privacy literature before deploying topic models.

\subsection{Energy and Carbon Footprint}
\label{sec:energy}

\textcite{strubell2019energy} showed that training large neural NLP
models can require energy equivalent to dozens of transatlantic flights.
While LDA and other classical topic models are computationally
inexpensive, the trend toward larger neural models and LLM-based
topic discovery reintroduces significant energy costs.
Researchers should report training and inference times, hardware
specifications, and estimated energy consumption alongside model
results, following the recommendations of \citeauthor{strubell2019energy}.
The environmental cost of research infrastructure should be a routine
part of the methods section, not an optional afterthought.

% =======================================================================
\section{Research Gaps and Future Directions}
\label{sec:future}
% =======================================================================

Our review identifies several open problems that deserve sustained
attention from the research community.

\paragraph{Evaluation reform.}
The evidence reviewed in Section~\ref{sec:evaluation} makes a
compelling case that the field needs standardised, multi-dimensional
evaluation protocols that combine automated metrics with human
judgment and downstream task performance.
The current state is one of fragmentation: coherence metrics vary
across papers (NPMI, $C_{v}$, UMass, UCI), reference corpora differ,
preprocessing pipelines are inconsistently reported, and human
evaluation is rare.
The OCTIS framework \parencite{terragni2021octis} is a step in this
direction, but it covers only a subset of existing models and metrics.
A community-maintained evaluation benchmark---analogous to GLUE in
sentence-level NLP---would substantially advance the field, provided
it was designed with domain coverage, language diversity, and
task-oriented metrics in mind.
Such a benchmark should include stability metrics (e.g.\ Jaccard
similarity between topics across random seeds) as first-class citizens,
not optional additions.

\paragraph{Principled hyperparameter selection.}
How to choose $K$ (the number of topics) remains an unsolved problem
despite decades of attention.
Bayesian approaches (HDP, DP mixtures) automate $K$ selection in
principle, but their performance is sensitive to hyperparameter
settings that are themselves difficult to tune.
Neural topic models add further hyperparameters (encoder depth, latent
dimensionality, learning rate, warm-up schedule) whose interaction
with topic quality is poorly understood.
Practical guidance for practitioners on how to choose and validate
$K$ and related hyperparameters is sparse and often contradictory.
A systematic, empirically grounded guide---analogous to Bayesian
optimisation guidelines in other areas of NLP---would have high
practical impact.

\paragraph{Stability as a first-class property.}
Following \textcite{koltcov2014lda}, stability should be routinely
measured and reported.
Models that produce high-coherence topics on one random seed but
qualitatively different topics on another are not reliable analytical
tools.
This is especially pressing for neural topic models, where weight
initialisation, data shuffling, and GPU non-determinism all contribute
to run-to-run variation.
Ensemble methods that aggregate topics across multiple runs---
collapsing near-duplicate topics and averaging document-topic
representations---are underexplored and could significantly improve
the reliability of topic-model outputs in practice.

\paragraph{Integration of probabilistic structure with LLMs.}
Current LLM-based topic discovery methods are largely pipeline-based.
A unified model that combines the LLM's world knowledge and language
generation capabilities with a principled probabilistic framework
would enable the best of both paradigms: coherent, readable topics
with well-calibrated uncertainty and document-level probability
estimates.
One promising research direction is to use an LLM's conditional
generation as the observation model within a variational
autoencoder framework, replacing the standard linear decoder with
a frozen or fine-tuned LLM.
This would preserve the probabilistic structure needed for downstream
use while leveraging the LLM's semantic richness.

\paragraph{Cross-lingual and low-resource settings.}
Despite progress in cross-lingual topic models
\parencite{bianchi2021crosslingual,mimno2009polylingual}, low-resource
languages with limited pre-training data remain inadequately served.
Methods that require less pre-training while still providing
cross-lingual alignment are needed, as are evaluation benchmarks
for cross-lingual topic quality that go beyond English-centred metrics.

\paragraph{Long-form and structured documents.}
Most topic models assume short-to-medium documents.
For books, legal documents, and multi-section scientific papers,
hierarchical or segment-aware topic models are more appropriate but
less well studied.
The interaction between discourse structure and topic assignment is a
rich area for future work: a legal brief, for instance, may contain
an argument section, a factual section, and a precedent section, each
with different topical emphases that a flat topic model will blend
together.
Incorporating structural priors---either from metadata or from
discourse parsing---into topic models remains an open problem.

\paragraph{Privacy-preserving topic modeling.}
As topic models are increasingly applied to sensitive text
(medical records, private communications, legal documents),
the development of differentially private inference algorithms
for topic models deserves more attention than it currently receives.
Existing approaches for differentially private LDA sacrifice either
accuracy or computational efficiency to a degree that limits
practical adoption; more efficient mechanisms are needed.

% =======================================================================
\section{Conclusion}
\label{sec:conclusion}
% =======================================================================

Topic modeling has matured from a single landmark paper---LDA---into
a broad and technically sophisticated field.
The trajectory from LSA through pLSA to Bayesian topic models and
then to neural variational autoencoders, embedding-based methods,
and LLM-assisted discovery reflects genuine scientific progress:
each generation addressed real limitations of its predecessors.
LSA established the statistical tractability of latent semantic
decomposition; pLSA added a principled probabilistic generative story;
LDA resolved the train/test disconnect with a Bayesian prior; neural
models scaled the framework to high-throughput GPU computation; and
LLM-based methods brought the richness of world knowledge into the
topic-discovery loop.

Yet our critical review reveals a persistent and troubling asymmetry:
the pace of model innovation has consistently outrun the development
of reliable evaluation methods.
Coherence metrics that the field treats as standards are poorly
calibrated against human judgment; stability is underreported;
extrinsic benchmarks are underused; and the most recent generation
of LLM-based methods imports the additional problems of cost,
reproducibility, and lack of probabilistic grounding.

The warnings from \textcite{chang2009reading}---that perplexity
anticorrelates with interpretability---were absorbed into the field's
collective knowledge, but the lesson was not generalised: the community
adopted automatic coherence metrics as replacements for perplexity
without subjecting those metrics to equal scrutiny.
The subsequent work of \textcite{hoyle2021automated} and
\textcite{doogan2021twaddle} shows that automatic coherence metrics are
themselves flawed proxies, subject to the same category of criticism
that was levelled at perplexity: they measure something, but not quite
what we want.
The lesson to draw is not that evaluation is impossible but that no
single metric is sufficient.

Progress on evaluation reform is not merely a methodological luxury.
Without it, the empirical claims that drive the field's development
rest on a shaky foundation.
We encourage researchers to adopt multi-dimensional evaluation
protocols, report stability across random seeds, include human
evaluation even at small scale, and use standardised benchmarking
frameworks such as OCTIS \parencite{terragni2021octis}.

Practitioners choosing a topic model for a specific task should be
guided by three questions.
First, does the use case require a probabilistic document
representation (e.g.\ for downstream modelling, retrieval, or
uncertainty quantification)?
If so, LLM-based approaches are currently unsuitable.
Second, is the corpus large enough to support the model's parameter
count reliably?
Neural models offer flexibility that comes at the cost of requiring
larger corpora to achieve stable results; on small corpora, LDA
remains the most reliable choice.
Third, is the evaluation criterion aligned with the downstream goal?
A model chosen to maximise $C_{v}$ coherence may not be the model
that best serves a classification or retrieval task, and vice versa.

The application landscape---from political science
\parencite{bail2016combining} to scientometrics
\parencite{griffiths2004finding} to public health---demonstrates that
topic models, when used carefully, provide genuine analytical value.
That value is best secured not by adopting the newest method
uncritically but by matching method to task, evaluation metric to
downstream goal, and model complexity to the evidence available in
the data.
The critical lens applied in this review is offered in that spirit:
not to diminish the achievements of the field, but to identify where
stronger foundations are needed and to provide a roadmap for building
them.

% =======================================================================
\printbibliography
% =======================================================================

\end{document}
